\chapter{The logarithmic cost of a Euclidean catenoidal neck}\label{sec:catenoidal-neck-models}
This section develops the explicit Euclidean calculation in \cite{KMN}, Section 2.1, Proposition 2.1 and Remark 2.2, pp.~4--6 of the supplied manuscript. It explains the scale in Remark \ref{rem:catenoid-estimate}. It does not construct the manifold sweepout of that paper's Theorem 2.4, and it makes no conclusion about multiplicity in an unspecified min-max class.

\begin{proposition}[Two rotational catenoids with the same circular boundary]\label{prop:catenoid-boundary-branches}
Fix $R>0$ and $h>0$. A rotational catenoid of waist radius $a>0$, centered at height zero, meets the planes $z=\pm h$ in circles of radius $R$ exactly when
\[
 R=a\cosh(h/a).
\]
There is a unique $t_*>0$ satisfying $t_*\tanh t_*=1$. Set $q_*=t_*/\cosh t_*$. If $0<h/R<q_*$, there are exactly two such waist radii. At equality there is exactly one, and above it there is none. The smaller waist corresponds to the larger solution $t=h/a>t_*$. The same threshold parameter occurs in the free-boundary catenoid of Example \ref{ex:critical-catenoid}; the present assertion counts rotational solutions for prescribed boundary circles.
\end{proposition}
\begin{proof}
The catenoid parametrization in Example \ref{ex:explicit-catenoid} has radius $a\cosh v$ at height $av$, giving the equation. On putting $t=h/a$, it becomes $h/R=t/\cosh t$. The derivative of the right side is $(1-t\tanh t)/\cosh t$. The function $t\tanh t$ is strictly increasing from zero to infinity, since its derivative is $\tanh t+t/\cosh^2t>0$. Thus the ratio increases strictly to a unique maximum and then decreases strictly to zero; its limit at zero is also zero. The number of roots follows by the intermediate value theorem and strict monotonicity. Since $a=h/t$, the larger root has the smaller waist. This argument counts rotational solutions and does not assert uniqueness among all surfaces spanning the two circles, or prove their stability indices.
\end{proof}

\begin{proposition}[Exact excess over two disks]\label{prop:catenoid-exact-disk-excess}
For a truncated catenoid with waist $a$ and parameter range $|v|\leq T$, put $h=aT$ and $R=a\cosh T$. Its area $A$ and excess over the two radius-$R$ disks are
\[
 A=2\pi a^2(T+\sinh T\cosh T),\qquad
 A-2\pi R^2=2\pi a^2\left(T-\frac12-\frac12e^{-2T}\right).
\]
In particular the excess is positive for all sufficiently large $T$, although its ratio to $h^2$ tends to zero as $T\to\infty$.
\end{proposition}
\begin{proof}
The induced metric is $a^2\cosh^2v(dv^2+d\theta^2)$. Integrating the area density over a full angular period and $[-T,T]$ gives the first formula. For the second, subtract $2\pi a^2\cosh^2T$ and use
\[
 \sinh T\cosh T-\cosh^2T=-\cosh T e^{-T}
 =-\frac12(1+e^{-2T}).
\]
Finally divide the excess by $h^2=a^2T^2$. The resulting expression is $2\pi(T-1/2-e^{-2T}/2)/T^2$, positive eventually and tending to zero. No limiting interchange with an area-minimizing problem is involved.
\end{proof}

\begin{theorem}[The sharp leading logarithmic neck cost]\label{thm:catenoid-logarithmic-excess}
Fix $R>0$ and let $h\downarrow0$. Choose the smaller waist solution $a(h)$ of Proposition \ref{prop:catenoid-boundary-branches}. If $L=\log(R/h)$ and $T=h/a(h)$, then
\[
 \frac{T}{L}\longrightarrow1,\qquad
 T=L+\log L+\log2+o(1),\qquad
 \frac{\bigl(A(h)-2\pi R^2\bigr)L}{h^2}\longrightarrow2\pi.
\]
Consequently, for every $\varepsilon>0$ and all sufficiently small $h>0$,
\[
 A(h)\leq2\pi R^2+\frac{2\pi(1+\varepsilon)h^2}{\log(R/h)}.
\]
This is the dimensionless version of the leading estimate in \cite{KMN}, Proposition 2.1 and Remark 2.2; it concerns the explicit rotational neck.
\end{theorem}
\begin{proof}
The larger root $T$ tends to infinity: on the strictly decreasing branch of $t/\cosh t$, its value $h/R$ tends to zero, and no bounded sequence on that branch can have this limit. The boundary equation yields
\[
 L=\log\cosh T-\log T
 =T-\log2+\log(1+e^{-2T})-\log T.
\]
Divide by $T$ to obtain $L/T\to1$. Rearrangement and $\log T-\log L\to0$ then give the second expansion. By Proposition \ref{prop:catenoid-exact-disk-excess},
\[
 \frac{(A-2\pi R^2)L}{h^2}
 =2\pi\frac{L}{T}\left(1-\frac{1+e^{-2T}}{2T}\right)\longrightarrow2\pi.
\]
The asserted bound is the definition of this limit, with $L>0$ once $h<R$. The logarithm uses a ratio of lengths, so it is invariant under a common change of length units.
\end{proof}

\begin{example}[A cylindrical connector has a larger leading cost]\label{ex:cylindrical-neck-cost}
For $0<h<R$, remove a disk of radius $\rho\in[0,R]$ from each of the planar radius-$R$ disks at heights $\pm h$, and join the inner boundary circles by a cylinder of radius $\rho$. The resulting piecewise smooth comparison set has area
\[
 2\pi R^2+4\pi h\rho-2\pi\rho^2.
\]
The maximum along this explicit radius family is $2\pi R^2+2\pi h^2$, attained at $\rho=h$. For fixed $R$ and $h\downarrow0$, this excess is larger by an asymptotic factor $\log(R/h)$ than the thin catenoid's excess.
\end{example}
\begin{proof}
The two annuli contribute $2\pi(R^2-\rho^2)$ and the cylinder, of height $2h$, contributes $4\pi h\rho$. The excess is $2\pi h^2-2\pi(\rho-h)^2$, which has the stated maximum because $h$ lies in the allowed interval. The ratio assertion follows from Theorem \ref{thm:catenoid-logarithmic-excess}. At the joins this comparison set has corners; the displayed calculation alone neither supplies a smooth sweepout with controlled smoothing errors nor identifies a variational width. Those additional requirements belong to the deformation theorem, not this elementary area comparison.
\end{proof}

\begin{proposition}[The exact planar annulus capacity]\label{prop:planar-annulus-exact-energy}
Let $0<a<R$ and $\mathcal A=\{x\in\mathbb R^2:a<|x|<R\}$. Among smooth functions on the closed annulus with boundary values zero on the inner circle and one on the outer circle, the unique minimizer of $\int_{\mathcal A}|\nabla u|^2$ is
\[
 u_0(x)=\frac{\log(|x|/a)}{\log(R/a)},\qquad
 \int_{\mathcal A}|\nabla u_0|^2=\frac{2\pi}{\log(R/a)}.
\]
For every such function $u$, the exact excess identity is
\[
 \int_{\mathcal A}|\nabla u|^2
 =\frac{2\pi}{\log(R/a)}+\int_{\mathcal A}|\nabla(u-u_0)|^2.
\]
\end{proposition}
\begin{proof}
The radial Laplacian of $\log|x|$ is zero away from zero, so $u_0$ is harmonic. Its radial derivative is $1/(|x|\log(R/a))$; polar integration gives the stated energy. Write $u=u_0+v$, where $v$ vanishes on both boundary circles. Integration by parts gives $\int\nabla u_0\cdot\nabla v=0$. Expanding the square gives the identity. Equality forces $\nabla v=0$, and the boundary value then forces $v=0$. This is an exact minimization on a fixed planar annulus. Lemma \ref{lem:point-capacity-cutoffs} already uses the logarithmic profile to remove punctures; the present identity specifies the optimal constant and its equality case.
\end{proof}

\begin{proposition}[Quadratic sheet gain dominates a logarithmic error]\label{prop:unstable-sheet-logarithmic-comparison}
Let $\Sigma$ be a smooth closed embedded two-sided minimal surface in a smooth three-manifold, with chosen unit normal $\nu$. Suppose $\phi>0$ is smooth and its area second variation satisfies $Q(\phi,\phi)=-q<0$. For small $t$, set
\[
 \Sigma_t=\{\exp_p(t\phi(p)\nu(p)):p\in\Sigma\}.
\]
Then
\[
 |\Sigma_t|+|\Sigma_{-t}|=2|\Sigma|-qt^2+O(t^4).
\]
If a collection of comparison surfaces $S_{t,s}$ additionally satisfies, uniformly in its parameter $s$,
\[
 |S_{t,s}|\leq |\Sigma_t|+|\Sigma_{-t}|
       +\frac{Ct^2}{\log(1/|t|)}
\]
for some fixed $C$ and all sufficiently small $0<|t|<1$, then every surface in that collection has area at most $2|\Sigma|-(q/2)t^2$ once $t$ is small enough. The existence, topology and continuity of such a collection are additional hypotheses here, not conclusions of this estimate.
\end{proposition}
\begin{proof}
For small $t$, the normal exponential map is a tubular coordinate map on the compact embedded surface. The positivity of $\phi$ makes the two graphs disjoint when $t\ne0$. Their area function $A(t)$ is smooth, with $A'(0)=0$ by minimality and $A''(0)=-q$ by Theorem \ref{thm:second-variation}. Taylor expansion through the cubic term gives
$A(t)+A(-t)=2A(0)-qt^2+O(t^4)$, because the odd terms cancel. Choose $t$ small enough that the absolute remainder is at most $(q/4)t^2$ and $C/\log(1/|t|)\leq q/4$. Substituting into the assumed uniform comparison gives the result. This algebra explains the competition described in \cite{KMN}, p.~2; proving a uniformly controlled manifold deformation requires the separate geometric construction of its Section 2.2.
\end{proof}
